\documentclass[11pt]{article}

\usepackage[margin=25mm]{geometry}

\usepackage{amsmath,amssymb,fontspec,unicode-math}

\IfFontExistsTF{DejaVu Serif}{\setmainfont{DejaVu Serif}}{\setmainfont{Latin Modern Roman}}

\IfFontExistsTF{Latin Modern Math}{\setmathfont{Latin Modern Math}}{\setmathfont{DejaVu Math TeX Gyre}}

\usepackage{xurl,hyperref}

\setlength{\parindent}{0pt}\setlength{\parskip}{6pt}

\begin{document}

\title{HAIF-GS: Hierarchical and Induced Flow-Guided Gaussian Splatting for Dynamic Scene}

\author{ELVIA preliminary AI analysis}\date{21 October 2026}\maketitle

Preliminary assessment: Minor.\par

Independent scientific review remains pending. Author adoption has not been requested.

Current tables still contain scene-mean inconsistencies, mixed-evaluator rankings and overly broad superiority wording. Motion/displacement and directional-flow notation also require local clarification. These are Minor; temporal variance, cycle-consistency identifiability and binary-penalty observations are retained as scope clarifications, without alleging core-model failure.

\section*{Review scope}

Full original concern list and current proceedings methods, experimental tables and appendix read; current author README/license checked at the resolved head. Direct arithmetic and identity-motion arguments were checked without executing author or historical verifier code.

\section*{Problem 1: Several scene means exceed rounding bounds}

Source: Paper pp. 7, 8, 15, 16

Assessment: supported; impact: minor

\subsection*{Source claim (paraphrase)}

Displayed means aggregate the listed scene values.

\subsection*{Analysis}

The displayed entries give 3DGS PSNR \(20.86\) rather than \(20.29\); TiNeuVox PSNR \(32.785714\) rather than \(31.92\), and LPIPS \(0.049571\) rather than \(0.038\); GauFre PSNR \(23.542857\) rather than 23.64. Seven metric aggregates exceed rounding bounds. Times \(38\), \(26\), \(41\), \(21\), \(23\), \(35\), \(27\) average \(30.142857\) minutes, incompatible with \(29\) under ordinary rounding. Some discrepancies understate the proposed method.

\subsection*{Finding and scope}

The bounded local discrepancy or counterexample is supported; no change in published benchmark scores is quantified.

\subsection*{Proposed correction}

Reconcile the specific equation, reported quantity or source contract described above. Corrected full-experiment results remain unverified.

Proposed correction; independent verification and author adoption remain unestablished.

\section*{Problem 2: Different evaluation functions share a ranking column}

Source: Paper pp. 7

Assessment: supported; impact: minor

\subsection*{Source claim (paraphrase)}

Values in each ranked metric column are comparable evaluation quantities.

\subsection*{Analysis}

Table 1 discloses mixed SSIM/MS-SSIM and Alex/VGG LPIPS, yet ranks the resulting values in common columns. Disclosure does not make different functions numerically interchangeable; PSNR must be assessed separately.

\subsection*{Finding and scope}

The bounded local discrepancy or counterexample is supported; no change in published benchmark scores is quantified.

\subsection*{Proposed correction}

Reconcile the specific equation, reported quantity or source contract described above. Corrected full-experiment results remain unverified.

Proposed correction; independent verification and author adoption remain unestablished.

\section*{Problem 3: Universal superiority exceeds the table evidence}

Source: Paper pp. 8–9

Assessment: supported; impact: minor

\subsection*{Source claim (paraphrase)}

The method leads across all scenes and metrics.

\subsection*{Analysis}

Cup PSNR is \(24.72\) versus NeRF-DS \(24.91\); Basin is \(19.74\) versus HyperNeRF 20.41. Several same-evaluator LPIPS values also favor baselines. The justified conclusion is metric- and scene-specific.

\subsection*{Finding and scope}

The bounded local discrepancy or counterexample is supported; no change in published benchmark scores is quantified.

\subsection*{Proposed correction}

Reconcile the specific equation, reported quantity or source contract described above. Corrected full-experiment results remain unverified.

Proposed correction; independent verification and author adoption remain unestablished.

\section*{Problem 4: Identity motion returns position rather than displacement}

Source: Paper pp. 6

Assessment: supported; impact: minor

\subsection*{Source claim (paraphrase)}

Eq.5 returns displacement and combines consistent rotation representations.

\subsection*{Analysis}

With identity rotations, zero translations and normalized weights, Eq. (\(5\)) returns the original center \(\mu\). Its left side is displacement, which should instead be zero under identity motion. The equation therefore needs either a position interpretation or subtraction of \(\mu\). It also combines axis-angle, matrix and quaternion quantities without conversion. For example, opposite unit quaternions \(q\) and \(-q\) represent one rotation but average to zero unless signs are aligned.

\subsection*{Finding and scope}

The motion formula is locally inconsistent. No released core implementation establishes that the same notation error produced the paper's reconstruction results.

\subsection*{Proposed correction}

Reconcile the specific equation, reported quantity or source contract described above. Corrected full-experiment results remain unverified.

Proposed correction; independent verification and author adoption remain unestablished.

\section*{Problem 5: Flow heads and time steps are underspecified}

Source: Paper pp. 5, 7

Assessment: supported; impact: minor

\subsection*{Source claim (paraphrase)}

Eq.7 defines a three-dimensional directional flow at the stated time offset.

\subsection*{Analysis}

Eq.7 combines a three-dimensional flow with a bidirectional network output without selecting a head. Adjacent normalized frame times are not generally \(t\pm1\).

\subsection*{Finding and scope}

The bounded local discrepancy or counterexample is supported; no change in published benchmark scores is quantified.

\subsection*{Proposed correction}

Reconcile the specific equation, reported quantity or source contract described above. Corrected full-experiment results remain unverified.

Proposed correction; independent verification and author adoption remain unestablished.

\section*{Problem 6: Translation variance is scale-dependent}

Source: Paper pp. 6, 14

Assessment: not\_issue; impact: no\_issue\_found

\subsection*{Source claim (paraphrase)}

The temporal translation variance is a heuristic for anchor refinement.

\subsection*{Analysis}

Equation (6) averages squared translation deviations over \(16\) sampled times. Repeating translations \(0\) and \((1,0,0)\) eight times each gives variance \(0.25\); doubling coordinates gives variance \(1\), crossing the same \(0.75\) threshold. This correctly exposes length-squared units. The paper calls the quantity normalized variance but does not assert coordinate-scale invariance; normalization can mean division by the number of time samples. A common rigid translation or pure rotation further limits its interpretation as nonrigid complexity, without invalidating its declared heuristic use.

\subsection*{Finding and scope}

The variance formula is valid. Its scene-scale and motion-type dependence is a scope clarification, not a demonstrated author error or reconstruction failure.

\section*{Problem 7: The author release lacks the core pipeline}

Source: Paper pp. 14

Assessment: needs\_evidence; impact: needs\_evidence

\subsection*{Source claim (paraphrase)}

The public release and manuscript specify the complete reconstruction pipeline.

\subsection*{Analysis}

Training, deformation, rendering and evaluation implementations are absent. Network widths, anchor schedules and run manifests do not determine a unique reimplementation. This does not establish that private experiments failed to run.

\subsection*{Finding and scope}

This is a scope/provenance gap, not proof that the stronger empirical conclusion is false.

\subsection*{Proposed correction}

Reconcile the specific equation, reported quantity or source contract described above. Corrected full-experiment results remain unverified.

Proposed correction; independent verification and author adoption remain unestablished.

\section*{Problem 8: Cycle consistency does not certify physical scene flow}

Source: Paper pp. 5–7

Assessment: not\_issue; impact: no\_issue\_found

\subsection*{Source claim (paraphrase)}

Cycle consistency encourages temporally coherent internal flow for reconstruction.

\subsection*{Analysis}

Zero forward and backward fields have zero cycle loss, as do arbitrary mutually inverse maps. Therefore cycle loss alone does not identify physical scene flow. Sections 4.3–4.5 describe internal temporal guidance and encouragement of coherence, without claiming an if-and-only-if identification theorem. The full training objective also includes reconstruction, so the zero-flow example does not prove that the trained method collapses or fails.

\subsection*{Finding and scope}

This limits a stronger identifiability interpretation that the paper does not establish. The reported reconstruction ablations remain legitimate evidence for reconstruction usefulness.

\section*{Problem 9: Efficiency is a metric-specific tradeoff}

Source: Paper pp. 16

Assessment: needs\_evidence; impact: needs\_evidence

\subsection*{Source claim (paraphrase)}

The method has the stated general efficiency advantage.

\subsection*{Analysis}

Rendering throughput is about \(1.60\) times 4DGS, but memory is \(3.3\) versus \(1.1\) GB and recomputed training time \(30.14\) versus \(26\) minutes. Matched SC-GS costs are absent, so no universal resource advantage follows.

\subsection*{Finding and scope}

This is a scope/provenance gap, not proof that the stronger empirical conclusion is false.

\subsection*{Proposed correction}

Reconcile the specific equation, reported quantity or source contract described above. Corrected full-experiment results remain unverified.

Proposed correction; independent verification and author adoption remain unestablished.

\section*{Problem 10: The loss-weight scan does not explain the default}

Source: Paper pp. 7, 17

Assessment: needs\_evidence; impact: needs\_evidence

\subsection*{Source claim (paraphrase)}

The loss-weight scan supports the reported default selection.

\subsection*{Analysis}

Table 12 reports PSNR \(25.57\) at weight \(1.0\), above final \(24.63\), while default \(0.2\) is absent from that scan. Different protocols or randomness could explain this; \(25.57\) must not be silently replaced by 24.57.

\subsection*{Finding and scope}

This is a scope/provenance gap, not proof that the stronger empirical conclusion is false.

\subsection*{Proposed correction}

Reconcile the specific equation, reported quantity or source contract described above. Corrected full-experiment results remain unverified.

Proposed correction; independent verification and author adoption remain unestablished.

\section*{Problem 11: The ratio denominator is missing}

Source: Paper pp. 16

Assessment: needs\_evidence; impact: needs\_evidence

\subsection*{Source claim (paraphrase)}

The displayed ratio has a reproducible denominator.

\subsection*{Analysis}

If Table 10 means anchors/Gaussians, Sieve gives \(3.7/114\simeq0.03246\) rather than 0.287. Another denominator is possible. Missing definition is supported; a numerical-error claim is conditional.

\subsection*{Finding and scope}

This is a scope/provenance gap, not proof that the stronger empirical conclusion is false.

\subsection*{Proposed correction}

Reconcile the specific equation, reported quantity or source contract described above. Corrected full-experiment results remain unverified.

Proposed correction; independent verification and author adoption remain unestablished.

\section*{Problem 12: The full weight interval does not emphasize the center}

Source: Paper pp. 6

Assessment: supported; impact: minor

\subsection*{Source claim (paraphrase)}

The stated weight interval emphasizes the central time.

\subsection*{Analysis}

For weights \((\lambda,1-2\lambda,\lambda)\), the middle weight exceeds an endpoint only if \(1-2\lambda>\lambda\), or \(\lambda<1/3\). The full stated interval \(\lambda<1/2\) is too broad. The actual setting \(\lambda=0.25\) satisfies the narrower condition.

\subsection*{Finding and scope}

The bounded local discrepancy or counterexample is supported; no change in published benchmark scores is quantified.

\subsection*{Proposed correction}

Reconcile the specific equation, reported quantity or source contract described above. Corrected full-experiment results remain unverified.

Proposed correction; independent verification and author adoption remain unestablished.

\section*{Problem 13: The binary term is not Shannon entropy}

Source: Paper pp. 7

Assessment: not\_issue; impact: no\_issue\_found

\subsection*{Source claim (paraphrase)}

The binary confidence penalty encourages endpoints and is combined with a sparsity penalty.

\subsection*{Analysis}

The penalty \(\alpha(1-\alpha)\) is Bernoulli variance, or a binary Gini/Tsallis-type impurity, rather than Shannon entropy. It is maximal at \(1/2\) and zero at either endpoint, exactly supporting the stated binary-selection explanation. The combined regularizer \(0.5\alpha+0.2\alpha(1-\alpha)\) has derivative \(0.7-0.4\alpha>0\) on \([0,1]\); its preference for zero is consistent with its declared sparsity purpose, while reconstruction supplies the competing incentive. The paper does not identify this function specifically as Shannon entropy.

\subsection*{Finding and scope}

The defined penalties have the stated endpoint/sparsity roles. Neither their naming nor this derivative proves collapse of the full model.

\section*{Problem 14: Experimental uncertainty is not inapplicable}

Source: Paper pp. 19–25

Assessment: needs\_evidence; impact: needs\_evidence

\subsection*{Source claim (paraphrase)}

Statistical reporting is inapplicable to these experiments.

\subsection*{Analysis}

Stochastic training experiments remain applicable to the statistical-reporting checklist. A displayed \(0.04\) dB gain does not by itself establish a stable cross-run advantage.

\subsection*{Finding and scope}

This is a scope/provenance gap, not proof that the stronger empirical conclusion is false.

\subsection*{Proposed correction}

Reconcile the specific equation, reported quantity or source contract described above. Corrected full-experiment results remain unverified.

Proposed correction; independent verification and author adoption remain unestablished.

\section*{Limits}

Independent scientific adjudication remains pending.\par

No author code, model training or full benchmark evaluation was executed in this revision.\par

Version-specific source defects do not establish which code produced the published experiments.\par

Corrected benchmark scores, empirical effect sizes and author adoption remain unverified unless explicitly stated.\par

No new author-code execution, training, empirical replication or contact.\par

Historical verification outputs were inspected but not rerun.\par

Missing run provenance and unresolved empirical interpretations remain explicit.\par

\section*{Source identity}

NeurIPS 2025 current proceedings PDF (checked 21 October 2026)

\url{https://proceedings.neurips.cc/paper_files/paper/2025/file/b690b88de1cd9694e356b021bc643ba1-Paper-Conference.pdf}

Reviewed PDF SHA-256: \nolinkurl{e5dfdfbde1b3017b8d06ebb946aaec7e63a0fc0a22ac9d8f03a44ed8e58dba68}

This current primary is the same research work as the predecessor, identified by matching title, authors and technical content, on the stated current publication channel. Its exact inspected bytes are identified by SHA-256; the earlier immutable source identity remains in the predecessor record. The download timestamp is recorded from this fresh primary-source retrieval.

\end{document}
